%% notes-tex/database/database-expansion-200/sections/rule.tex
%% [[experiments.database.expansion-200]]
%% https://github.com/Mjvolk3/torchcell/tree/main/notes-tex/database/database-expansion-200/sections/rule.tex
%% SOURCE: hand-authored prose; every number quoted here comes from
%% experiments/database/scripts/build_candidate_datasets_table.py

\section{The selection rule\secstatus{todo}}\label{sec:rule}

\subsection{Sequence reconstruction is the gate, not a preference\secstatus{todo}}

A phenotype is only trainable against a genotype that can be written down. Every row in
Table~\ref{tab:final} therefore names a \term{sequence basis}: the route by which the
total genomic content of one strain would be reconstructed. The eleven bases in use are

\begin{itemize}[leftmargin=1.2em,itemsep=1pt,topsep=2pt]
  \item \term{S288C-KO} -- the reference genome minus one cataloged ORF (BY4741, BY4742,
    BY4743). Thirteen of the sixty in Table~\ref{tab:final} sit here, second only to
    reference-only.
  \item \term{S288C-KO/het} -- a heterozygous diploid deletion, so a dosage edit rather
    than a null.
  \item \term{S288C+guide} -- the genome is unedited and the perturbation is a designed
    cassette plus a guide target. Regulatory, not sequence-level.
  \item \term{S288C+designed-edit} -- a designed single-nucleotide or small-indel library.
    Designed, and not verified per strain.
  \item \term{S288C+tag} -- the reference plus a designed tag, degron or promoter cassette.
  \item \term{S288C+ORF-plasmid} -- the reference plus a barcoded ORF on a known plasmid.
  \item \term{S288C+reporter-locus} -- the reference plus a designed regulatory sequence
    driving a reporter, which is what every massively parallel reporter assay perturbs.
  \item \term{isolate-WGS} -- a published assembly or variant call set per isolate.
  \item \term{segregant-WGS} -- genotype calls or whole-genome sequence per progeny.
  \item \term{engineered-chassis} -- a named production strain plus heterologous cassettes.
  \item \term{reference-only} -- wild type, where the environment carries the perturbation.
\end{itemize}

What fails the gate is a strain whose genome is unknown even in principle. Genome shuffling
and adaptive laboratory evolution without resequencing both land there: the selected
progeny differ from their parents at unmapped positions, so there is no genotype to map a
phenotype from. Those are dropped in Table~\ref{tab:excluded} rather than ranked low, since
no amount of phenotypic value substitutes for a missing genotype.

The \term{S288C+designed-edit} and \term{S288C+guide} bases carry a weaker claim than the
others and the distinction is worth holding. A designed variant library states what each
member was built to be, not what it turned out to be, so its genotypes are designed rather
than realized. Ingesting one commits to storing the design and a verification status, not a
verified sequence.

\subsection{Band comes before tier, and why\secstatus{todo}}\label{sec:band}

Scale is the right default ordering and it is the wrong first question when a specific
campaign is being planned. The largest row in this table, de Boer 2020 at $10^8$ promoter
sequences, tells a model of state across genomes nothing; 112 sequenced strains with a
transcriptome, a proteome and a phosphoproteome each tell it what the pairing looks like.
Each row therefore carries a \term{band}, applied before tier and before scale:

\begin{itemize}[leftmargin=1.2em,itemsep=1pt,topsep=2pt]
  \item \term{state across genomes}, 8 rows. A transcriptome or proteome measured on
    strains whose genomes differ by sequence: isolates, segregants, a recombinant panel or
    an interspecies hybrid. The owner's focus for this pass, set on 2026-10-19, and the
    axis Sec.~\ref{sec:state} shows the served graph is thinnest on. It replaces the
    previous pass's \emph{next build} band, whose three rows (Muenzner 2024, Peter 2018,
    De Chiara 2022) are served.
  \item \term{perturb-seq}, 23 rows. A CRISPR interference or activation library whose
    per-guide matrix is released, a barcoded genotype set read out per cell, a designed
    single-gene perturbation crossed with a transcriptome, or the strain and noise
    measurements a pooled per-cell campaign depends on. Seven rows of the previous pass's
    thirty are served (Boocock 2025, Hale 2024, Puddu 2019, Hackett 2020, Hu 2007,
    N'Guessan 2025, Jackson 2020).
  \item \term{molecular layers}, 14 rows. Two or more of transcript, translation,
    protein, phosphosite, metabolite, flux and turnover, joinable on one axis, either
    inside one study or against a served dataset. The join is per \emph{strain} when
    both layers are measured on the same genotypes, and per \emph{gene} when one layer is
    a single genome-wide coefficient such as a half-life. The eight state rows were in
    this band last pass and are promoted out of it.
  \item \term{scale}, 109 rows. Everything else, ordered as before.
\end{itemize}

Within a band the order is tier, then measurements. The cost is visible and is not hidden:
the three bands total 45 rows, so they fill ranks 1 to 45 and the first five scale rows
complete the fifty: de Boer 2020, Lee 2014, Nguyen Ba 2022, Galardini 2019 and Dutta 2026.
The ten extra are all scale. de Boer 2020 and Lee 2014 are the largest things in the
ranking by instance count and sit at 46 and 47, which is the deliberate consequence of
ordering by what a row combines with rather than by how many numbers it holds. A reader
who wants scale back reads rows 46 to 60 first.

\subsection{Tier, and what rows are ranked on\secstatus{todo}}

Tier is assigned from stated bars rather than judgment:

\begin{itemize}[leftmargin=1.2em,itemsep=1pt,topsep=2pt]
  \item \textbf{Tier 1} clears all three of $\geq$$10^3$ genotypes, $\geq$$10^4$ instances,
    and a clean sequence basis, and lands in one of the nine classes.
  \item \textbf{Tier 2} clears two of the three, or is the only dataset covering a
    requested class, or directly de-risks the Perturb-seq proposal.
  \item \textbf{Tier 3} is a modality or reference backbone the substrate needs without
    itself carrying a large genotype axis.
  \item \textbf{Tier 4} is real but small, coarsely labeled, or expensive to extract.
\end{itemize}

The bars apply to what a row contains rather than to its reputation. Turco 2023 has the
widest environment axis in the pool at 7,536 and is tier 2, because it is a
meta-aggregation: its instance count re-serves primary screens, several of them already
served, and the net-new fraction is unknown until it is de-duplicated against the rest of
this table.

Within a tier, rows sort by \term{measurements}, meaning instances times phenotype
dimensionality. Instances alone is the wrong statistic and gets the order visibly wrong:
Albert 2018 is 1,012 instances, which ranks it below a mid-sized fitness screen, but each
instance is a transcriptome, so it is roughly $5.8\times10^6$ numbers. This is the same
normalization the supported-datasets table already applies through its gzip-signal column.

\subsection{Where the lines fall, and what they cost\secstatus{todo}}

The long-run cut is 200 additions, $270-70$, and the pool holds 154 rows, so every row in
the pool is inside the cut and there is no reserve. The 46-row shortfall is not padded;
Sec.~\ref{sec:state-gaps} names where the next discovery pass should look.

The two nearer lines are what a build week is planned against. Wave 1 is rows 1--50 and
wave 2 rows 51--70; a wave-2 row is promoted the moment a wave-1 row turns out to have no
recoverable per-strain data, which has already happened once to Lee 2014. Wave 1 is the
eight state rows, 23 Perturb-seq rows, 14 molecular-layer rows and five scale rows; wave 2
is scale throughout, so the two waves differ in kind and not only in rank.
Table~\ref{tab:counts} gives the composition by class and by band.

%% GENERATED FILE -- do not hand-edit.
%% SOURCE: experiments/database/scripts/build_candidate_datasets_table.py
\begin{table}[H]\centering
\small
\caption[]{Candidates by class and by band, split at the two wave lines.
\emph{Genotypes} and \emph{Instances} sum the per-row axes over all waves; rows with
no count contribute nothing, so both totals are lower bounds.}
\label{tab:counts}
\begin{tabular}{@{}l r r r r r@{}}
\toprule
 & Wave 1 & Wave 2 & Reserve & Genotypes & Instances \\
\midrule
Natural variation & 14 & 4 & 14 & 1,225,745 & $3.1\times 10^{6}$ \\
Tolerance / robustness & 3 & 4 & 15 & 122,906 & $5.3\times 10^{7}$ \\
CRISPR library screen & 8 & 0 & 3 & 328,794 & $3.9\times 10^{5}$ \\
Expression / single cell & 15 & 1 & 3 & 529 & $1.0\times 10^{3}$ \\
Metabolite / precursor & 5 & 2 & 7 & 40,631 & $1.3\times 10^{5}$ \\
Modality / backbone & 4 & 4 & 29 & 785,185 & $1.3\times 10^{6}$ \\
Regulatory DNA & 1 & 1 & 5 & 100,509,217 & $1.0\times 10^{8}$ \\
Deep mutational scan & 0 & 4 & 5 & 811,600 & $8.2\times 10^{5}$ \\
Combinatorial genome & 0 & 0 & 3 & 400 & 600 \\
\midrule
\multicolumn{6}{@{}l}{\emph{The same rows, by band}}\\
state across genomes & 8 & 0 & 0 & 3,000 & $3.1\times 10^{3}$ \\
perturb-seq & 23 & 0 & 0 & 327,436 & $4.0\times 10^{5}$ \\
molecular layers & 14 & 0 & 0 & 5,500 & $5.7\times 10^{3}$ \\
scale & 5 & 20 & 84 & 103,489,071 & $1.6\times 10^{8}$ \\
\midrule
Total & 50 & 20 & 84 & 103,825,007 & $1.6\times 10^{8}$ \\
\bottomrule
\end{tabular}
\end{table}


A dataset named explicitly in the scoping request is pinned into the recommended set even
if it ranks below the cut, and the pin is reported rather than folded into the score, so
the set stays auditable. With the whole pool inside the cut no pin can bind, so none is
in force; the two requested rows, Blank 2005 and Trikka 2015, rank on their own figures.

\subsection{What moved since the previous pass\secstatus{todo}}\label{sec:moves}

Table~\ref{tab:swaps} lists every position change against the committed ranking of the
previous document, bounded to rows inside either pass's working set. Thirteen rows left
the list because KG 4.0 serves them, each with the dataset class it is served as. The
eight state rows climbed from ranks 34 to 112 into the first eight. Everything else in the
table moved because rows above it left or were promoted, and says so.

%% GENERATED FILE -- do not hand-edit.
%% SOURCE: experiments/database/scripts/build_candidate_datasets_table.py
\begingroup
\footnotesize
\setlength{\tabcolsep}{4pt}
\renewcommand{\arraystretch}{1.15}
\begin{longtable}{@{}L{46mm} r@{\hspace{6pt}} r@{\hspace{6pt}} L{104mm}@{}}
\caption[]{Rank changes between the previous pass and this one, for every row in
either pass's first 70, every row added, and every row that left.
\emph{Was} is recomputed under the previous rule, tier then measurements, over the
rows that pass held. A dash means the row is new or has gone.}
\label{tab:swaps}\\
\toprule
\textbf{Dataset} & \textbf{Was} & \textbf{Now} & \textbf{Reason} \\
\midrule
\endfirsthead
\toprule
\textbf{Dataset} & \textbf{Was} & \textbf{Now} & \textbf{Reason} \\
\midrule
\endhead
\bottomrule
\endfoot
\textbf{Albert 2018 (eQTL in 1,012 segregants)} & 34 & 1 & State across genomes, first: bulk RNA-seq of 1,012 sequenced BY x RM segregants, the largest genotype-indexed transcriptome in yeast, and the bulk partner of the served Boocock 2025, whose pooled panel stores 393 of these same segregants. \\
\addlinespace[5pt]
\textbf{Jakobson 2025 (genome-to-proteome map)} & 35 & 2 & State across genomes: protein abundance across the same sequenced segregant panel Albert 2018 profiles by transcriptome, 800 F6 segregants, so protein and transcript loci are measurable on one genotype set. \\
\addlinespace[5pt]
\textbf{Teyssonniere 2024 (species-wide proteome against transcriptome)} & 37 & 3 & State across genomes: proteomes for 942 of the 1,011 isolates whose transcriptomes the served Caudal 2024 holds and whose proteomes the served Muenzner 2024 holds for 796, so the de-duplication against Muenzner is the first task and the pairing with Caudal is one to one. \\
\addlinespace[5pt]
\textbf{Grossbach 2022 (BY x RM transcriptome, proteome and phosphoproteome)} & 38 & 4 & State across genomes: transcript, protein and phosphosite on one set of 112 sequenced BY x RM strains, the only row where all three layers are observed within a genotype. \\
\addlinespace[5pt]
\textbf{Smith 2008 (gene-environment interaction eQTL)} & 112 & 5 & State across genomes: mRNA abundance of 109 BY x RM segregants in glucose and in ethanol, the one segregant transcriptome with a condition axis, so genotype by environment is a within-panel contrast. \\
\addlinespace[5pt]
\textbf{Skelly 2013 (expression variation across isolates)} & 46 & 6 & State across genomes: transcriptome and proteome on the same 22 sequenced isolates, the cheap and the expensive layer paired within one experiment on natural genomes. \\
\addlinespace[5pt]
\textbf{McManus 2014 (ribosome profiling, allele-specific)} & 48 & 7 & State across genomes: ribosome occupancy beside mRNA abundance for S. cerevisiae, S. paradoxus and their F1 hybrid, so translation is read allele by allele across two genomes; the paradoxus arm is out of species and is stored apart or dropped. \\
\addlinespace[5pt]
\textbf{Foss 2007 (BY x RM segregant proteome)} & 53 & 8 & State across genomes: the first segregant proteome on the BY x RM cross and the independent replicate of Grossbach 2022 fifteen years and two platforms earlier. \\
\addlinespace[5pt]
\textbf{Dong 2021 (MAGIC + SAM biosensor)} & 11 & 9 & A genome-wide CRISPR activation, interference and deletion library sorted on a metabolite biosensor, so the label is product concentration rather than fitness. The pattern a sorted Perturb-seq would reuse. \\
\addlinespace[5pt]
\textbf{Momen-Roknabadi 2020 (inducible CRISPRi library)} & 12 & 10 & A genome-scale inducible CRISPR interference library with a released per-guide matrix. Induction control is what lets a knockdown be applied after a cell is captured rather than during outgrowth. \\
\addlinespace[5pt]
\textbf{McGlincy 2021 (genome-scale CRISPRi library)} & 13 & 11 & The second independently designed genome-scale yeast CRISPR interference library with a released per-guide matrix, so guide design and gene effect can be separated. \\
\addlinespace[5pt]
\textbf{Bao 2018 (CHAnGE single-nucleotide library)} & 14 & 12 & A guide-indexed genome-wide library whose edits sit below the open reading frame, so the perturbation is an allele rather than a null. \\
\addlinespace[5pt]
\textbf{Crook 2016 (tunable RNAi, isobutanol + 1-butanol)} & 15 & 13 & A dose-graded knockdown axis. Graded rather than binary perturbation is what turns a per-cell readout into a dose-response curve. \\
\addlinespace[5pt]
\textbf{Roy 2018 (multiplexed precision editing)} & 16 & 14 & Multiplexed designed edits per cell, which is combinatorial input in the sense the Perturb-seq specification requires. \\
\addlinespace[5pt]
\textbf{Newman 2006 (protein noise, GFP library)} & 17 & 15 & Per-cell protein abundance distributions across a genome-scale tagged collection. Expression noise per gene is what sets how large a perturbation effect has to be before a per-cell readout can resolve it. \\
\addlinespace[5pt]
\textbf{Mukherjee 2021 (CRISPRi essential genes x acetic acid)} & 18 & 16 & Knockdown reaches the essential genes a deletion collection cannot hold, which is a third of the genome a deletion-based campaign would miss. \\
\addlinespace[5pt]
\textbf{Lian 2017 (CRISPR-AID, beta-carotene)} & 19 & 17 & Three perturbation modalities in one cell, activation, interference and deletion, which is combinatorial input by modality rather than by gene. \\
\addlinespace[5pt]
\textbf{Hughes 2000 (compendium of expression profiles)} & 20 & 18 & The original designed-perturbation expression compendium, and the independent measurement that makes cross-laboratory replication of the deletion transcriptome testable against the built Kemmeren 2014. \\
\addlinespace[5pt]
\textbf{Lenstra 2011 (chromatin regulator deletion expression)} & 21 & 19 & The same design over chromatin regulators rather than sequence-specific factors, so the two together cover both halves of transcriptional control by deletion. \\
\addlinespace[5pt]
\textbf{Sun 2013 (mRNA synthesis and decay rates across deletion strains)} & 22 & 20 & A designed single-gene perturbation crossed with a transcriptome, and the only one that resolves the level into a synthesis rate and a decay rate. Rates are what a per-cell snapshot cannot recover on its own. \\
\addlinespace[5pt]
\textbf{Jariani 2020 (yeast scRNA-seq through lag phase)} & 23 & 21 & Per-cell states through a carbon-source shift, so the readout is the distribution across an unsynchronized population rather than its mean. \\
\addlinespace[5pt]
\textbf{Guo 2018 (CRISPR-Cas9 tiling of essential genes)} & 24 & 22 & Guide tiling produces a graded allelic series across one gene, so guide position rather than gene identity carries the effect. \\
\addlinespace[5pt]
\textbf{Urbonaite 2021 (yeastDrop-Seq under drug treatment)} & 25 & 23 & Single-drug arms against their combination with a per-cell readout, the environmental analog of a combinatorial genetic perturbation. \\
\addlinespace[5pt]
\textbf{Nadal-Ribelles 2019 (sensitive yeast scRNA-seq)} & 26 & 24 & The high-sensitivity yeast protocol, strand and isoform aware. Sensitivity per cell is the constraint on how small a perturbation effect a campaign can resolve. \\
\addlinespace[5pt]
\textbf{Gasch 2017 (single-cell stress heterogeneity)} & 27 & 25 & Separates intrinsic from extrinsic per-cell variation in an isogenic population under stress, which is the null a perturbation effect is measured against. \\
\addlinespace[5pt]
\textbf{Jaffe 2019 (multiplexed CRISPR interference epistasis)} & 28 & 26 & Two knockdowns in the same cell with a measured interaction, which is the only combinatorial CRISPR interference design located in yeast. \\
\addlinespace[5pt]
\textbf{Jackson 2023 (wild-type scRNA-seq time course for RNA kinetics)} & 29 & 27 & The platform companion of Jackson 2020: 173,361 cells of one strain, which is what sets the per-cell variance a perturbation has to exceed. \\
\addlinespace[5pt]
\textbf{Brettner 2024 (ultra-high-throughput yeast scRNA-seq)} & 30 & 28 & Combinatorial barcoding in yeast at 96 to 384 multiplexed genotypes or environments per run. The throughput route by which a campaign becomes affordable. \\
\addlinespace[5pt]
\textbf{Mulleder 2012 (prototrophic deletion collection)} & 31 & 29 & The prototrophic deletion collection. Jackson 2020 used a prototrophic background because auxotrophy blocks minimal and nitrogen-limited media, so this is the strain resource a campaign in defined media needs. \\
\addlinespace[5pt]
\textbf{Su 2023 (single-cell transcriptomes under four stresses)} & 32 & 30 & Full-length rather than three-prime tag counting, so isoform-level readout is on the table for a campaign that needs it. \\
\addlinespace[5pt]
\textbf{Wang 2022 (single-cell transcriptomes across replicative aging)} & 33 & 31 & The only yeast single-cell set with age as the condition axis, which is a covariate any pooled campaign carries whether or not it measures it. \\
\addlinespace[5pt]
\textbf{Aulakh 2025 (genome-scale ionome)} & 36 & 32 & A metabolic readout on the whole deletion collection, which is the genotype axis the supported expression compendium already covers. \\
\addlinespace[5pt]
\textbf{Leutert 2023 (phosphoproteome x 101 conditions)} & 39 & 33 & Enzyme regulation acts faster than transcription, so a phosphosite layer over 101 conditions is what explains flux changes an expression table cannot. \\
\addlinespace[5pt]
\textbf{Brauer 2008 (growth-rate-controlled chemostat transcriptome)} & 40 & 34 & Expression at controlled growth rate under each nutrient limitation, which is the condition axis Boer 2010 and Hackett 2016 measure metabolites and flux on. \\
\addlinespace[5pt]
\textbf{Hackett 2016 (SIMMER multi-omic flux)} & 41 & 35 & Flux, metabolite and transcript measured in one study under the same nutrient limitations, so the join is internal rather than across papers. \\
\addlinespace[5pt]
\textbf{Zhu 2014 (kinase / phosphatase lipidomics)} & 42 & 36 & A lipidome over 129 signaling deletions, most of which carry an expression profile in the supported compendium. \\
\addlinespace[5pt]
\textbf{Gerke 2017 (urea-cycle mQTL)} & 43 & 37 & Metabolite quantitative trait loci on a segregant cross that also has an expression map, so a metabolite locus can be read through its transcript. \\
\addlinespace[5pt]
\textbf{Ambroset 2014 (metabolite QTL)} & 44 & 38 & A 74-metabolite panel on a sequenced segregant panel, the widest metabolite vector available on a recombinant genotype axis. \\
\addlinespace[5pt]
\textbf{Blank 2005 (13C metabolic flux)} & 45 & 39 & The only row measuring flux rather than a concentration, on deletion mutants that also have a transcriptome in the supported set. \\
\addlinespace[5pt]
\textbf{Airoldi 2016 (nitrogen-limited steady-state and dynamic transcriptome)} & 47 & 40 & Expression under the exact nitrogen-limited conditions that Jackson 2020 profiles single cells in, so the deletion effect separates from the medium effect. \\
\addlinespace[5pt]
\textbf{Martin-Perez 2017 (protein half-lives)} & 49 & 41 & Protein turnover per gene, the second of those two terms, spanning three orders of magnitude and therefore not replaceable by a global constant. \\
\addlinespace[5pt]
\textbf{Boer 2010 (metabolome across nutrient limitations)} & 50 & 42 & The metabolite half of the chemostat nutrient-limitation series whose transcriptome half is Brauer 2008, same laboratory and same conditions. \\
\addlinespace[5pt]
\textbf{Tengolics 2024 (domestication metabolome)} & 51 & 43 & Metabolite levels across sequenced isolates, the panel the supported Caudal 2024 transcriptomes and Muenzner 2024 proteomes also sit on. \\
\addlinespace[5pt]
\textbf{Eder 2020 (flux QTL)} & 52 & 44 & Flux mapped to segregant genotypes, which is the flux counterpart of the expression and protein maps on the same panel type. \\
\addlinespace[5pt]
\textbf{Yu 2021 (proteome and metabolome under nitrogen limitation)} & 54 & 45 & Both layers measured in one study under nitrogen limitation, so it calibrates the protein-to-metabolite step that cross-study joins assume. \\
\addlinespace[5pt]
\textbf{de Boer 2020 (100M random promoters)} & 55 & 46 & Unchanged criteria; moved by 9 as rows around it were promoted into a band. \\
\addlinespace[5pt]
\textbf{Lee 2014 (HIP-HOP fitness signatures)} & 57 & 47 & Unchanged criteria; moved by 10 as rows around it were promoted into a band. \\
\addlinespace[5pt]
\textbf{Nguyen Ba 2022 (barcoded bulk QTL, 100k segregants)} & 58 & 48 & Unchanged criteria; moved by 10 as rows around it were promoted into a band. \\
\addlinespace[5pt]
\textbf{Galardini 2019 (four backgrounds x 38 conditions)} & 59 & 49 & Unchanged criteria; moved by 10 as rows around it were promoted into a band. \\
\addlinespace[5pt]
\textbf{Dutta 2026 (barcoded natural-isolate chemical response)} & 61 & 50 & Unchanged criteria; moved by 11 as rows around it were promoted into a band. \\
\addlinespace[5pt]
\textbf{Li 2016 (tRNA fitness landscape)} & 62 & 51 & Unchanged criteria; moved by 11 as rows around it were promoted into a band. \\
\addlinespace[5pt]
\textbf{Domingo 2018 (tRNA double-mutant landscape)} & 63 & 52 & Unchanged criteria; moved by 11 as rows around it were promoted into a band. \\
\addlinespace[5pt]
\textbf{Zhang 2021 (YETI titratable overexpression)} & 64 & 53 & Unchanged criteria; moved by 11 as rows around it were promoted into a band. \\
\addlinespace[5pt]
\textbf{Peeters 2021 (fermentation-trait QTL atlas)} & 65 & 54 & Unchanged criteria; moved by 11 as rows around it were promoted into a band. \\
\addlinespace[5pt]
\textbf{Van Leeuwen 2024 (short-chain organic acids)} & 66 & 55 & Unchanged criteria; moved by 11 as rows around it were promoted into a band. \\
\addlinespace[5pt]
\textbf{Chica 2026 (AutoDRY autophagy screen)} & 67 & 56 & Unchanged criteria; moved by 11 as rows around it were promoted into a band. \\
\addlinespace[5pt]
\textbf{Sopko 2006 (genome-scale overexpression toxicity)} & 68 & 57 & Unchanged criteria; moved by 11 as rows around it were promoted into a band. \\
\addlinespace[5pt]
\textbf{Liu 2021 (tryptophan / isobutanol tolerance)} & 69 & 58 & Unchanged criteria; moved by 11 as rows around it were promoted into a band. \\
\addlinespace[5pt]
\textbf{Kuroda 2019 (isobutanol-specific tolerance)} & 70 & 59 & Unchanged criteria; moved by 11 as rows around it were promoted into a band. \\
\addlinespace[5pt]
\textbf{Turco 2023 (Yeast Phenome)} & 71 & 60 & Unchanged criteria; moved by 11 as rows around it were promoted into a band. \\
\addlinespace[5pt]
\textbf{Rossi 2021 (ChIP-exo protein architecture)} & 72 & 61 & Unchanged criteria; moved by 11 as rows around it were promoted into a band. \\
\addlinespace[5pt]
\textbf{Bloom 2025 (global epistasis in a yeast cross)} & 74 & 62 & Unchanged criteria; moved by 12 as rows around it were promoted into a band. \\
\addlinespace[5pt]
\textbf{Gasch 2000 (environmental stress response)} & 75 & 63 & Unchanged criteria; moved by 12 as rows around it were promoted into a band. \\
\addlinespace[5pt]
\textbf{Decourty 2021 (RNA-metabolism genetic interactions)} & 76 & 64 & Unchanged criteria; moved by 12 as rows around it were promoted into a band. \\
\addlinespace[5pt]
\textbf{Cuperus 2017 (random 5' UTR library)} & 77 & 65 & Unchanged criteria; moved by 12 as rows around it were promoted into a band. \\
\addlinespace[5pt]
\textbf{Faure 2022 (doubledeepPCA)} & 78 & 66 & Unchanged criteria; moved by 12 as rows around it were promoted into a band. \\
\addlinespace[5pt]
\textbf{Chong 2015 / CYCLoPs (single-cell proteome dynamics)} & 79 & 67 & Unchanged criteria; moved by 12 as rows around it were promoted into a band. \\
\addlinespace[5pt]
\textbf{Teyssonniere 2024 (species-wide trait survey)} & 80 & 68 & Unchanged criteria; moved by 12 as rows around it were promoted into a band. \\
\addlinespace[5pt]
\textbf{Diss 2018 (Fos-Jun interaction double mutants)} & 81 & 69 & Unchanged criteria; moved by 12 as rows around it were promoted into a band. \\
\addlinespace[5pt]
\textbf{Warringer 2011 (trait variation across the SGRP panel)} & 82 & 70 & Unchanged criteria; moved by 12 as rows around it were promoted into a band. \\
\addlinespace[5pt]
\textbf{Piotrowski 2017 (MOSAIC diagnostic panel)} & 73 & -- & Served since the previous pass in KG 4.0 (release 2026.10.17-4295f7ca) as EnvChemgenPiotrowski2017Dataset; it is now a join partner rather than a candidate. \\
\addlinespace[5pt]
\textbf{Hale 2024 (CRISPRi x natural variation)} & 5 & -- & Served since the previous pass in KG 4.0 (release 2026.10.17-4295f7ca) as CrispriSegregantHale2024Dataset; it is now a join partner rather than a candidate. \\
\addlinespace[5pt]
\textbf{Parsons 2006 (bioactive-compound profiling)} & 60 & -- & Served since the previous pass in KG 4.0 (release 2026.10.17-4295f7ca) as EnvChemgenParsons2006Dataset; it is now a join partner rather than a candidate. \\
\addlinespace[5pt]
\textbf{Peter 2018 (1,011 isolate genomes + phenome)} & 2 & -- & Served since the previous pass in KG 4.0 (release 2026.10.17-4295f7ca) as IsolateGrowthPeter2018Dataset; it is now a join partner rather than a candidate. \\
\addlinespace[5pt]
\textbf{Muenzner 2024 (natural-isolate proteome)} & 1 & -- & Served since the previous pass in KG 4.0 (release 2026.10.17-4295f7ca) as two datasets (ProteomeMuenzner2024Dataset and ProteinTurnoverMuenzner2024Dataset); it is now a join partner rather than a candidate. \\
\addlinespace[5pt]
\textbf{Puddu 2019 (WGS of the deletion collection)} & 7 & -- & Served since the previous pass in KG 4.0 (release 2026.10.17-4295f7ca) as GenomeArchitecturePuddu2019Dataset; it is now a join partner rather than a candidate. \\
\addlinespace[5pt]
\textbf{N'Guessan 2025 (segregant scRNA-seq eQTL)} & 8 & -- & Served since the previous pass in KG 4.0 (release 2026.10.17-4295f7ca) as Nguessan2025Dataset; it is now a join partner rather than a candidate. \\
\addlinespace[5pt]
\textbf{Vaishnav 2022 (regulatory DNA fitness landscape)} & 56 & -- & Served since the previous pass in KG 4.0 (release 2026.10.17-4295f7ca) as PooledPromoterExpressionVaishnav2022Dataset; it is now a join partner rather than a candidate. \\
\addlinespace[5pt]
\textbf{Boocock 2025 (single-cell eQTL mapping)} & 4 & -- & Served since the previous pass in KG 4.0 (release 2026.10.17-4295f7ca) as Boocock2025Dataset, 79,516 records over four segregant panels; it is now a join partner rather than a candidate. \\
\addlinespace[5pt]
\textbf{De Chiara 2022 (domestication and the life cycle)} & 3 & -- & Served since the previous pass in KG 4.0 (release 2026.10.17-4295f7ca) as three datasets (DoublingTimeDeChiara2022Dataset, CellYieldDeChiara2022Dataset, LifeCycleDeChiara2022Dataset); it is now a join partner rather than a candidate. \\
\addlinespace[5pt]
\textbf{Hu 2007 (TF deletion expression compendium)} & 10 & -- & Served since the previous pass in KG 4.0 (release 2026.10.17-4295f7ca) as MicroarrayScoreHu2007Dataset; it is now a join partner rather than a candidate. \\
\addlinespace[5pt]
\textbf{Jackson 2020 (TF-deletion single-cell atlas)} & 6 & -- & Served since the previous pass in KG 4.0 (release 2026.10.17-4295f7ca) as TfDeletionPseudobulkJackson2020Dataset; it is now a join partner rather than a candidate. \\
\addlinespace[5pt]
\textbf{Hackett 2020 (IDEA inducible-TF transcriptome time series)} & 9 & -- & Served since the previous pass in KG 4.0 (release 2026.10.17-4295f7ca) as InductionTimecourseHackett2020Dataset; it is now a join partner rather than a candidate. \\
\addlinespace[5pt]
\end{longtable}
\endgroup

